\subsection*{الف}
سری فوریه مثلثی را به صورت
\[
f(t)=\frac{a_0}{2}+\sum_{n=1}^{\infty}\big(a_n\cos nt+b_n\sin nt\big)
\]
در نظر می‌گیریم. از روابط داده‌شده داریم
\[
a_n=\frac{\pi n}{n^2+1}= \frac{\pi}{n}-\frac{\pi}{n(n^2+1)}
\quad\Longrightarrow\quad
n a_n \xrightarrow[n\to\infty]{}\pi\neq 0,
\qquad
b_n=\frac{\pi}{n^2+\pi^2+n\pi}=O\!\left(\frac{1}{n^2}\right).
\]
اگر $f$ روی یک دوره پیوسته و (قطعه‌به‌قطعه) مشتق‌پذیر باشد، با یک بار جزءبه‌جزء کردن داریم
\[
a_n=-\frac{1}{\pi n}\int f'(t)\sin(nt)\,dt,\qquad
b_n=\frac{1}{\pi n}\int f'(t)\cos(nt)\,dt
\]
و طبق لم ریمن--لبگ، $\int f'(t)e^{jnt}dt\to 0$، لذا باید $n a_n\to 0$ و $n b_n\to 0$ شود؛
اما این با $n a_n\to\pi$ در تناقض است. پس $f$ نمی‌تواند همزمان پیوسته و (قطعه‌به‌قطعه) مشتق‌پذیر باشد؛
بنابراین در هر دوره دست‌کم یک ناپیوستگی (پرش) دارد و مشتق کلاسیک آن روی کل دوره وجود ندارد
(به زبان سیگنال‌ها، $f'(t)$ شامل مؤلفه‌های ضربه‌ای در نقاط پرش است).

\subsection*{ب}
تابع
\[
f(t)=t-t^2,\qquad 0<t<1
\]
را اگر صرفاً با دوره‌ی $1$ تکرار کنیم، چون $f(0)=f(1)=0$ پیوستگی برقرار است ولی
\[
f'(t)=1-2t,\qquad f'(0^+)=1,\ \ f'(1^-)= -1
\]
و در اتصالِ انتهای دوره، مشتق پرش دارد؛ لذا ضرایب سری فوریه از مرتبه‌ی $\frac{1}{n^2}$ کاهش می‌یابند (طبق راهنما).

برای رسیدن به مرتبه‌ی $\frac{1}{n^3}$ کافی است سیگنال متناوب حاصل روی هر دوره $C^1$ باشد (یعنی هم خود تابع و هم مشتقش در نقاط اتصالِ دوره پیوسته شوند).
یک گسترش ساده با دوره‌ی $2$:
\[
g(t)=
\begin{cases}
t-t^2, & 0<t<1,\\
t^2-3t+2=(t-1)(t-2), & 1<t<2,
\end{cases}
\qquad g(t+2)=g(t).
\]
در این صورت
\[
g(0)=g(1)=g(2)=0,
\qquad
g'(t)=
\begin{cases}
1-2t, & 0<t<1,\\
2t-3, & 1<t<2,
\end{cases}
\]
و داریم
\[
g'(1^-)= -1 = g'(1^+),\qquad g'(0^+)=1=g'(2^-),
\]
پس $g$ روی کل دوره $C^1$ است، ولی $g''$ در $t=1$ (و به تبعِ تناوب در $t=2k$) پرش دارد؛
بنابراین (با دو بار جزءبه‌جزء کردن) ضرایب سری فوریه‌ی $g$ از مرتبه‌ی
\[
a_n,\ b_n = O\!\left(\frac{1}{n^3}\right)
\]
به صفر همگرا می‌شوند.